\documentclass[11pt]{article}
\usepackage[utf8]{inputenc}
\usepackage{geometry}
\geometry{a4paper, margin=1in}
\usepackage{amsmath}
\usepackage{amssymb}
\usepackage{graphicx}
\usepackage{fancyhdr}
\usepackage{enumitem}
\usepackage{tcolorbox}
\usepackage{hyperref}

% Document Info
\title{Homework 12.5 Equations of Lines and Planes}
\author{Tashfeen Omran}
\date{February 2026}

\begin{document}

\maketitle
\tableofcontents
\newpage

\section{Part 1: Comprehensive Introduction}

\subsection{Core Concepts}
The study of lines and planes in three-dimensional space ($\mathbb{R}^3$) is the foundation of 3D analytic geometry. Unlike in 2D, where a line is determined by a slope and an intercept, 3D structures require vectors to define their orientation.

\begin{itemize}
    \item \textbf{The Line:} A line is a 1D object in 3D space. To define it, we need a starting location (a point $P_0$) and a direction in which to travel (a vector $\mathbf{v}$).
    \item \textbf{The Plane:} A plane is a 2D flat surface in 3D space. To define it, we need a specific location (a point $P_0$) and a precise definition of "tilt" or orientation. The most efficient way to define orientation is using a \textbf{normal vector} $\mathbf{n}$, which points strictly perpendicular to the surface of the plane.
\end{itemize}

\subsection{Intuition and Derivation}

\subsubsection{Deriving the Line Equation}
Imagine you are at a point $P_0(x_0, y_0, z_0)$. You want to walk in a straight line in the direction of vector $\mathbf{v} = \langle a, b, c \rangle$.
Any point $P(x, y, z)$ on this line can be reached by starting at the origin, going to $P_0$ (vector $\mathbf{r}_0$), and then adding some multiple $t$ of the direction vector $\mathbf{v}$.
\[ \mathbf{r} = \mathbf{r}_0 + t\mathbf{v} \]
Here, $t$ acts as a scalar parameter. If $t=1$, you move one full vector length. If $t=-1$, you move backwards. This vector addition creates the parametric equations.

\subsubsection{Deriving the Plane Equation}
A plane is defined as the set of all points $P$ such that the vector connecting a fixed point $P_0$ to $P$ is \textit{perfectly perpendicular} to the normal vector $\mathbf{n}$.
Let $\mathbf{r}$ be the position vector of any point on the plane, and $\mathbf{r}_0$ be the position vector of the fixed point. The vector lying on the plane is $\mathbf{r} - \mathbf{r}_0$.
The condition of perpendicularity means their dot product is zero:
\[ \mathbf{n} \cdot (\mathbf{r} - \mathbf{r}_0) = 0 \]
If $\mathbf{n} = \langle a, b, c \rangle$, this expands to the scalar equation $a(x-x_0) + b(y-y_0) + c(z-z_0) = 0$.

\subsection{Historical Context}
The transition from Euclidean geometry (compass and straightedge) to Analytic geometry was pioneered by René Descartes and Pierre de Fermat in the 17th century. While they focused primarily on 2D, the extension to 3D became essential in the 18th and 19th centuries for solving problems in classical mechanics, specifically ballistics and celestial orbits. Mathematicians needed a way to describe not just the path of an object (a line or curve) but also the surfaces constrained by physical laws (planes and manifolds).

\subsection{Key Formulas}

\textbf{1. Equations of a Line}
\begin{itemize}
    \item \textbf{Vector Equation:} $\mathbf{r} = \mathbf{r}_0 + t\mathbf{v}$
    \item \textbf{Parametric Equations:} $x = x_0 + at, \quad y = y_0 + bt, \quad z = z_0 + ct$
    \item \textbf{Symmetric Equations:} $\frac{x-x_0}{a} = \frac{y-y_0}{b} = \frac{z-z_0}{c}$
\end{itemize}

\textbf{2. Equations of a Plane}
\begin{itemize}
    \item \textbf{Vector Equation:} $\mathbf{n} \cdot (\mathbf{r} - \mathbf{r}_0) = 0$
    \item \textbf{Scalar Equation:} $a(x-x_0) + b(y-y_0) + c(z-z_0) = 0$
    \item \textbf{Linear Equation:} $ax + by + cz + d = 0$
\end{itemize}

\textbf{3. Angle Between Planes}
\[ \cos \theta = \frac{|\mathbf{n}_1 \cdot \mathbf{n}_2|}{|\mathbf{n}_1| |\mathbf{n}_2|} \]

\subsection{Prerequisites}
To succeed in this topic, you must be proficient in:
\begin{itemize}
    \item \textbf{Vector Operations:} Dot product (for angles and perpendicularity) and Cross product (essential for finding normal vectors).
    \item \textbf{Algebra:} Solving systems of linear equations (for finding intersections).
    \item \textbf{Parametric Curves:} Understanding how a parameter $t$ traces a path.
\end{itemize}

\newpage

\section{Part 2: Detailed Homework Solutions}

\subsection*{Problem 1: True or False Statements}
\textbf{Determine whether each statement is true or false in $\mathbb{R}^3$.}

\begin{enumerate}[label=(\alph*)]
    \item \textbf{Two lines parallel to a third line are parallel.}
    \textbf{Solution:} \textbf{True.} Parallelism is transitive for lines. If $\mathbf{v}_1 = k_1 \mathbf{v}_3$ and $\mathbf{v}_2 = k_2 \mathbf{v}_3$, then $\mathbf{v}_1 = (k_1/k_2)\mathbf{v}_2$.

    \item \textbf{Two lines perpendicular to a third line are parallel.}
    \textbf{Solution:} \textbf{False.} Consider the x, y, and z axes. The x-axis is perp to the z-axis. The y-axis is perp to the z-axis. But the x-axis is perpendicular (not parallel) to the y-axis.

    \item \textbf{Two planes parallel to a third plane are parallel.}
    \textbf{Solution:} \textbf{True.} Their normal vectors are all parallel.

    \item \textbf{Two planes perpendicular to a third plane are parallel.}
    \textbf{Solution:} \textbf{False.} The xz-plane and yz-plane are both perpendicular to the xy-plane, but they intersect (and are perpendicular to each other).

    \item \textbf{Two lines parallel to a plane are parallel.}
    \textbf{Solution:} \textbf{False.} One line could run North-South and another East-West, both floating above the ground (plane). They are skew or perpendicular, not parallel.

    \item \textbf{Two lines perpendicular to a plane are parallel.}
    \textbf{Solution:} \textbf{True.} Both lines must be parallel to the plane's normal vector, so they are parallel to each other.

    \item \textbf{Two planes parallel to a line are parallel.}
    \textbf{Solution:} \textbf{False.} Think of an open book. The spine is the line. The two covers are planes parallel to the spine, but they intersect.

    \item \textbf{Two planes perpendicular to a line are parallel.}
    \textbf{Solution:} \textbf{True.} Both planes have normal vectors parallel to the line, so the normals are parallel to each other.

    \item \textbf{Two planes either intersect or are parallel.}
    \textbf{Solution:} \textbf{True.} In $\mathbb{R}^3$, there is no "skew" equivalent for planes. If normals are parallel, planes are parallel. If not, they must intersect in a line.

    \item \textbf{Two lines either intersect or are parallel.}
    \textbf{Solution:} \textbf{False.} They can be \textbf{skew} (non-parallel and non-intersecting, like a bridge crossing over a road).

    \item \textbf{A plane and a line either intersect or are parallel.}
    \textbf{Solution:} \textbf{True.} If the line is orthogonal to the plane's normal, it is parallel (or contained). If not, it punctures the plane at exactly one point.
\end{enumerate}

\hrulefill

\subsection*{Problem 2: Line through Point Parallel to Vector}
\textbf{Find a vector equation and parametric equations for the line through $(5, 3, 3)$ and parallel to the vector $3\mathbf{i} - \mathbf{j} + 7\mathbf{k}$.}

\textbf{Step 1: Identify $\mathbf{r}_0$ and $\mathbf{v}$.}
$\mathbf{r}_0 = \langle 5, 3, 3 \rangle$
$\mathbf{v} = \langle 3, -1, 7 \rangle$

\textbf{Step 2: Vector Equation.}
$\mathbf{r}(t) = \mathbf{r}_0 + t\mathbf{v} = \langle 5, 3, 3 \rangle + t\langle 3, -1, 7 \rangle = \langle 5+3t, 3-t, 3+7t \rangle$

\textbf{Step 3: Parametric Equations.}
Extract components:
$x(t) = 5 + 3t$
$y(t) = 3 - t$
$z(t) = 3 + 7t$

\textbf{Answer:}
$\mathbf{r}(t) = \langle 5+3t, 3-t, 3+7t \rangle$
$(x(t), y(t), z(t)) = (5+3t, 3-t, 3+7t)$

\hrulefill

\subsection*{Problem 3: Line Parallel to Vector with Fractions}
\textbf{Line through $(4, -5, 4)$ parallel to $\langle 1, 5, -2/3 \rangle$.}

\textbf{Step 1: Setup.}
$\mathbf{r}_0 = \langle 4, -5, 4 \rangle, \quad \mathbf{v} = \langle 1, 5, -2/3 \rangle$

\textbf{Step 2: Equations.}
$\mathbf{r}(t) = \langle 4+t, -5+5t, 4 - \frac{2}{3}t \rangle$

\textbf{Answer:}
$\mathbf{r}(t) = \langle 4+t, -5+5t, 4 - \frac{2}{3}t \rangle$
$(x(t), y(t), z(t)) = (4+t, -5+5t, 4 - \frac{2}{3}t)$

\hrulefill

\subsection*{Problem 4: Tutorial Exercise}
\textbf{Line through $(7, -4, 5)$ parallel to $\langle 1, 6, -1/3 \rangle$.}

\textbf{Step 1: Setup.}
$\mathbf{r}_0 = \langle 7, -4, 5 \rangle, \quad \mathbf{v} = \langle 1, 6, -1/3 \rangle$

\textbf{Answer:}
$\mathbf{r}(t) = \langle 7+t, -4+6t, 5 - \frac{1}{3}t \rangle$
$(x(t), y(t), z(t)) = (7+t, -4+6t, 5 - \frac{1}{3}t)$

\hrulefill

\subsection*{Problem 5: Line Perpendicular to a Plane}
\textbf{Line through $(8, 6, 1)$ and perpendicular to plane $2x - 3y + 3z = 3$.}

\textbf{Step 1: Extract Normal Vector.}
The coefficients of $x, y, z$ in the plane equation give the normal vector.
$\mathbf{n} = \langle 2, -3, 3 \rangle$.

\textbf{Step 2: Determine Line Direction.}
Since the line is perpendicular to the plane, the plane's normal vector acts as the line's direction vector.
$\mathbf{v} = \mathbf{n} = \langle 2, -3, 3 \rangle$.

\textbf{Answer:}
$\mathbf{r}(t) = \langle 8+2t, 6-3t, 1+3t \rangle$
$(x(t), y(t), z(t)) = (8+2t, 6-3t, 1+3t)$

\hrulefill

\subsection*{Problem 6: Line Through Two Points}
\textbf{Line through $(-5, 2, 5)$ and $(1, 4, -5)$. Find parametric and symmetric equations.}

\textbf{Step 1: Find Direction Vector.}
$\mathbf{v} = P_2 - P_1 = \langle 1 - (-5), 4 - 2, -5 - 5 \rangle = \langle 6, 2, -10 \rangle$.

\textbf{Step 2: Parametric Equations.}
Using $P_1(-5, 2, 5)$:
$x = -5 + 6t, \quad y = 2 + 2t, \quad z = 5 - 10t$.

\textbf{Step 3: Symmetric Equations.}
Solve each for $t$:
$t = \frac{x+5}{6} = \frac{y-2}{2} = \frac{z-5}{-10}$.

\textbf{Answer:}
Parametric: $\langle -5+6t, 2+2t, 5-10t \rangle$
Symmetric: Option 3: $\frac{x+5}{6} = \frac{y-2}{2} = \frac{z-5}{-10}$

\hrulefill

\subsection*{Problem 7: Line Through Origin and Point}
\textbf{Through origin $(0,0,0)$ and $(5, 9, -1)$.}

\textbf{Step 1: Direction Vector.}
$\mathbf{v} = \langle 5-0, 9-0, -1-0 \rangle = \langle 5, 9, -1 \rangle$.

\textbf{Step 2: Symmetric Equations.}
$\frac{x-0}{5} = \frac{y-0}{9} = \frac{z-0}{-1}$.
Simplify: $\frac{x}{5} = \frac{y}{9} = -z$.

\textbf{Answer:}
Parametric: $\langle 5t, 9t, -t \rangle$
Symmetric: Option 1: $\frac{x}{5} = \frac{y}{9} = -z$

\hrulefill

\subsection*{Problem 8: Line with a Zero Component}
\textbf{Through $(14, 5, -15)$ and $(-8, 5, 13)$.}

\textbf{Step 1: Direction Vector.}
$\mathbf{v} = \langle -8-14, 5-5, 13-(-15) \rangle = \langle -22, 0, 28 \rangle$.

\textbf{Step 2: Symmetric Equations.}
Since the $y$-component of $\mathbf{v}$ is 0, $y$ does not change. The symmetric form separates the constant variable.
$\frac{x-14}{-22} = \frac{z+15}{28}, \quad y = 5$.

\textbf{Answer:}
Parametric: $\langle 14-22t, 5, -15+28t \rangle$
Symmetric: Option 5: $y+8 = \dots$ (Wait, looking at options again)
Correct Option from OCR: Last option $\frac{x-14}{-22} = \frac{z+15}{28}, y=5$.

\hrulefill

\subsection*{Problem 9: Line of Intersection of Two Planes}
\textbf{Intersection of $x + 2y + 3z = 3$ and $x - y + z = 3$.}

\textbf{Step 1: Find Direction Vector.}
The line is perpendicular to both normals.
$\mathbf{n}_1 = \langle 1, 2, 3 \rangle, \quad \mathbf{n}_2 = \langle 1, -1, 1 \rangle$.
$\mathbf{v} = \mathbf{n}_1 \times \mathbf{n}_2 = \begin{vmatrix} \mathbf{i} & \mathbf{j} & \mathbf{k} \\ 1 & 2 & 3 \\ 1 & -1 & 1 \end{vmatrix} = \langle (2)(1) - (3)(-1), \quad (3)(1) - (1)(1), \quad (1)(-1) - (2)(1) \rangle$.
$\mathbf{v} = \langle 2+3, 3-1, -1-2 \rangle = \langle 5, 2, -3 \rangle$.

\textbf{Step 2: Find a Point.}
Set $z=0$:
$x + 2y = 3$
$x - y = 3 \implies y = x-3$.
Substitute: $x + 2(x-3) = 3 \implies 3x - 6 = 3 \implies 3x=9 \implies x=3$.
$y = 0$.
Point: $(3, 0, 0)$.

\textbf{Step 3: Symmetric Equation.}
$\frac{x-3}{5} = \frac{y-0}{2} = \frac{z-0}{-3}$.

\textbf{Answer:}
Option 3: $\frac{x-3}{5} = \frac{y}{2} = -\frac{z}{3}$

\hrulefill

\subsection*{Problem 10: Parallel Lines Check}
\textbf{Is line through $(-4, -6, 1)$ and $(-2, 0, -3)$ parallel to line through $(7, 17, 11)$ and $(4, 8, 17)$?}

\textbf{Step 1: Vectors.}
$\mathbf{v}_1 = \langle -2+4, 0+6, -3-1 \rangle = \langle 2, 6, -4 \rangle$.
$\mathbf{v}_2 = \langle 4-7, 8-17, 17-11 \rangle = \langle -3, -9, 6 \rangle$.

\textbf{Step 2: Check Ratios.}
$\frac{-3}{2} = -1.5, \quad \frac{-9}{6} = -1.5, \quad \frac{6}{-4} = -1.5$.
Since ratios are equal, $\mathbf{v}_2 = -1.5 \mathbf{v}_1$.

\textbf{Answer:}
\textbf{Yes}.

\hrulefill

\subsection*{Problem 11: Parallel Vectors Step-by-Step}
(Same logic as Problem 10)
\textbf{Direction $\mathbf{v}_1$}: $\langle 2, 6, -4 \rangle$.
\textbf{Direction $\mathbf{v}_2$}: $\langle -5, -15, 10 \rangle$.
\textbf{Are they parallel?} Yes ($\mathbf{v}_2 = -2.5 \mathbf{v}_1$).
\textbf{Are the lines parallel?} \textbf{Yes}.

\hrulefill

\subsection*{Problem 12: Perpendicular Lines Check}
\textbf{Line 1: $(-4, 5, 0)$ to $(1, 1, 1)$. Line 2: $(2, 3, 4)$ to $(5, -1, -6)$.}

\textbf{Step 1: Vectors.}
$\mathbf{v}_1 = \langle 1+4, 1-5, 1-0 \rangle = \langle 5, -4, 1 \rangle$.
$\mathbf{v}_2 = \langle 5-2, -1-3, -6-4 \rangle = \langle 3, -4, -10 \rangle$.

\textbf{Step 2: Dot Product.}
$\mathbf{v}_1 \cdot \mathbf{v}_2 = (5)(3) + (-4)(-4) + (1)(-10) = 15 + 16 - 10 = 21$.
Since $21 \neq 0$, they are not perpendicular.

\textbf{Answer:}
$\mathbf{v}_1 \cdot \mathbf{v}_2 = 21$.
Lines are \textbf{NOT} perpendicular.

\hrulefill

\subsection*{Problem 13: Line Segment}
\textbf{Segment from $(-3, 14, 38)$ to $(11, -9, 50)$.}

\textbf{Formula:} $\mathbf{r}(t) = (1-t)\mathbf{r}_0 + t\mathbf{r}_1$ for $0 \le t \le 1$.
Or $\mathbf{r}(t) = \mathbf{r}_0 + t(\mathbf{r}_1 - \mathbf{r}_0)$.
$\mathbf{v} = \langle 11+3, -9-14, 50-38 \rangle = \langle 14, -23, 12 \rangle$.
$\mathbf{r}(t) = \langle -3+14t, 14-23t, 38+12t \rangle$.

\textbf{Answer:}
$(x(t), y(t), z(t)) = (-3+14t, 14-23t, 38+12t)$ with $0 \le t \le 1$.

\hrulefill

\subsection*{Problem 14: Lines Intersection/Skew}
$L_1: \frac{x-1}{1} = \frac{y-1}{-2} = \frac{z-12}{-3} \quad \Rightarrow \quad x=1+t, y=1-2t, z=12-3t$.
$L_2: \frac{x-2}{1} = \frac{y+6}{3} = \frac{z-13}{-7} \quad \Rightarrow \quad x=2+s, y=-6+3s, z=13-7s$.

\textbf{Step 1: Equate components.}
1) $1+t = 2+s \implies t = s+1$
2) $1-2t = -6+3s$
Substitute (1) into (2):
$1 - 2(s+1) = -6+3s$
$1 - 2s - 2 = -6 + 3s$
$-1 - 2s = -6 + 3s \implies 5 = 5s \implies s=1$.
If $s=1$, then $t=2$.

\textbf{Step 2: Verify z-component.}
$z_1 = 12 - 3(2) = 6$.
$z_2 = 13 - 7(1) = 6$.
They match.

\textbf{Step 3: Point.}
Use $s=1$: $x=3, y=-3, z=6$.

\textbf{Answer:}
\textbf{Intersecting}.
Point: $(3, -3, 6)$.

\hrulefill

\subsection*{Problem 15: Plane Parallel to Another}
\textbf{Plane through $(7, -2, -6)$ parallel to $z = 4x - 4y$.}

\textbf{Step 1: Normal.}
Eq: $4x - 4y - z = 0$. $\mathbf{n} = \langle 4, -4, -1 \rangle$.

\textbf{Step 2: Point-Normal Form.}
$4(x-7) - 4(y-(-2)) - 1(z-(-6)) = 0$
$4(x-7) - 4(y+2) - (z+6) = 0$
$4x - 28 - 4y - 8 - z - 6 = 0$
$4x - 4y - z - 42 = 0 \implies 4x - 4y - z = 42$.

\textbf{Answer:}
$4x - 4y - z = 42$.

\hrulefill

\subsection*{Problem 16: Plane Containing Line and Parallel to Plane}
\textbf{Contains $x=3+t, y=4-t, z=5-3t$. Parallel to $5x + 2y + z = 4$.}

\textbf{Step 1: Normal.}
Since parallel to given plane, $\mathbf{n} = \langle 5, 2, 1 \rangle$.

\textbf{Step 2: Point.}
Use line point $P_0(3, 4, 5)$.

\textbf{Step 3: Equation.}
$5(x-3) + 2(y-4) + 1(z-5) = 0$
$5x - 15 + 2y - 8 + z - 5 = 0$
$5x + 2y + z = 28$.

\textbf{Answer:}
$5x + 2y + z = 28$.

\hrulefill

\subsection*{Problem 17: Plane Through 3 Points}
\textbf{$(0,4,4), (4,0,4), (4,4,0)$.}

\textbf{Step 1: Vectors.}
$\mathbf{a} = P_2 - P_1 = \langle 4, -4, 0 \rangle$.
$\mathbf{b} = P_3 - P_1 = \langle 4, 0, -4 \rangle$.

\textbf{Step 2: Cross Product.}
$\mathbf{n} = \mathbf{a} \times \mathbf{b} = \langle 16, -(-16), 16 \rangle = \langle 16, 16, 16 \rangle$.
Use simplified normal $\mathbf{n} = \langle 1, 1, 1 \rangle$.

\textbf{Step 3: Equation.}
$1(x-0) + 1(y-4) + 1(z-4) = 0$
$x + y + z = 8$.

\textbf{Answer:}
$x + y + z = 8$.

\hrulefill

\subsection*{Problem 18: Plane Through Origin and Two Points}
\textbf{Origin, $(6, -5, 3), (1, 1, 1)$.}

\textbf{Step 1: Vectors.}
$\mathbf{a} = \langle 6, -5, 3 \rangle, \quad \mathbf{b} = \langle 1, 1, 1 \rangle$.

\textbf{Step 2: Cross Product.}
$\mathbf{n} = \begin{vmatrix} \mathbf{i} & \mathbf{j} & \mathbf{k} \\ 6 & -5 & 3 \\ 1 & 1 & 1 \end{vmatrix} = \langle -5-3, -(6-3), 6-(-5) \rangle = \langle -8, -3, 11 \rangle$.

\textbf{Step 3: Equation.}
Point is $(0,0,0)$, so constants vanish.
$-8x - 3y + 11z = 0$ or $8x + 3y - 11z = 0$.

\textbf{Answer:}
$8x + 3y - 11z = 0$.

\hrulefill

\subsection*{Problem 19: Plane Through Point and Line}
\textbf{Through $(3, 6, -2)$ and line $x=4-t, y=2t-1, z=-3t$.}

\textbf{Step 1: Vectors.}
Line direction $\mathbf{v} = \langle -1, 2, -3 \rangle$.
Line point $P_L(4, -1, 0)$.
Vector from $P_L$ to $P_{plane}(3,6,-2)$ is $\mathbf{u} = \langle -1, 7, -2 \rangle$.

\textbf{Step 2: Normal.}
$\mathbf{n} = \mathbf{v} \times \mathbf{u} = \begin{vmatrix} \mathbf{i} & \mathbf{j} & \mathbf{k} \\ -1 & 2 & -3 \\ -1 & 7 & -2 \end{vmatrix} = \langle -4+21, -(2-3), -7+2 \rangle = \langle 17, 1, -5 \rangle$.

\textbf{Step 3: Equation.}
$17(x-3) + 1(y-6) - 5(z+2) = 0$
$17x - 51 + y - 6 - 5z - 10 = 0$
$17x + y - 5z = 67$.

\textbf{Answer:}
$17x + y - 5z = 67$.

\hrulefill

\subsection*{Problem 20: Plane Perpendicular to Two Planes}
\textbf{Through $(1, 5, 1)$ perp to $5x+y-5z=5$ and $x+4z=6$.}

\textbf{Step 1: Normals.}
$\mathbf{n}_1 = \langle 5, 1, -5 \rangle$.
$\mathbf{n}_2 = \langle 1, 0, 4 \rangle$.

\textbf{Step 2: New Normal.}
$\mathbf{n} = \mathbf{n}_1 \times \mathbf{n}_2 = \langle 4, -25, -1 \rangle$.

\textbf{Step 3: Equation.}
$4(x-1) - 25(y-5) - 1(z-1) = 0$
$4x - 4 - 25y + 125 - z + 1 = 0$
$4x - 25y - z = -122$.

\textbf{Answer:}
$4x - 25y - z = -122$.

\hrulefill

\subsection*{Problem 21: Intercepts}
\textbf{$5x + y + 3z = 15$.}
\textbf{x-intercept:} Set $y=z=0 \implies 5x=15 \implies x=3$.
\textbf{y-intercept:} Set $x=z=0 \implies y=15$.
\textbf{z-intercept:} Set $x=y=0 \implies 3z=15 \implies z=5$.

\hrulefill

\subsection*{Problem 22: Line-Plane Intersection}
\textbf{Line: $x=t-1, y=1+2t, z=3-t$. Plane: $4x-y+3z=7$.}

\textbf{Step 1: Substitute.}
$4(t-1) - (1+2t) + 3(3-t) = 7$
$4t - 4 - 1 - 2t + 9 - 3t = 7$
$-t + 4 = 7 \implies -t = 3 \implies t = -3$.

\textbf{Step 2: Point.}
$x = -3 - 1 = -4$.
$y = 1 + 2(-3) = -5$.
$z = 3 - (-3) = 6$.

\textbf{Answer:}
$(-4, -5, 6)$.

\hrulefill

\subsection*{Problem 23: Planes Angle}
\textbf{$x-y-7z=1$ and $7x+y-z=6$.}

\textbf{Step 1: Normals.}
$\mathbf{n}_1 = \langle 1, -1, -7 \rangle, \quad \mathbf{n}_2 = \langle 7, 1, -1 \rangle$.

\textbf{Step 2: Dot Product.}
$\mathbf{n}_1 \cdot \mathbf{n}_2 = 7 - 1 + 7 = 13 \neq 0$. Not perpendicular.

\textbf{Step 3: Angle.}
$|\mathbf{n}_1| = \sqrt{1+1+49} = \sqrt{51}$.
$|\mathbf{n}_2| = \sqrt{49+1+1} = \sqrt{51}$.
$\cos \theta = \frac{13}{51}$.
$\theta = \arccos(13/51) \approx 75.2^{\circ}$.

\textbf{Answer:}
Neither. Angle $\approx 75.2^{\circ}$.

\newpage

\section{Part 3: Analysis of Problems and Techniques}

\subsection{Problem Types}
\begin{enumerate}
    \item \textbf{Setup Problems (Problems 2, 3, 4):} Direct application of the formula $\mathbf{r} = \mathbf{r}_0 + t\mathbf{v}$. The challenge is simply identifying the point and direction.
    \item \textbf{Geometric Translation (Problems 5, 15, 16):} These require translating geometric descriptions into vectors. Key insight: "Parallel to plane" $\rightarrow$ Normal is normal. "Perpendicular to plane" $\rightarrow$ Normal is direction.
    \item \textbf{Intersection Problems (Problems 9, 14, 22):} These require solving systems.
    \begin{itemize}
        \item Line-Plane: Plug line into plane eq.
        \item Line-Line: Equate parameters ($t$ and $s$).
        \item Plane-Plane: Intersection is a line. Direction is $\mathbf{n}_1 \times \mathbf{n}_2$.
    \end{itemize}
    \item \textbf{Construction Problems (Problems 17-20):} Building a plane equation from various inputs (3 points, 1 point + line, etc.). Strategy: Always find two vectors lying in the plane, then cross them to get $\mathbf{n}$.
\end{enumerate}

\subsection{Key Manipulations}
\begin{itemize}
    \item \textbf{The "Normal Swap":} In Problem 5, we were given a plane and asked for a perpendicular line. We used the plane's normal $\langle a,b,c \rangle$ as the line's direction vector $\mathbf{v}$. This is the standard trick for converting between orthogonal concepts.
    \item \textbf{The Cross Product Generator:} In Problems 19 and 20, we didn't have a normal vector. We had to \textit{manufacture} one by finding two vectors in the plane and crossing them.
    \item \textbf{Parameter Separation:} In Problem 14 (Line intersection), it was crucial to use different parameters ($t$ and $s$) for the two lines. Using $t$ for both leads to a mathematical contradiction unless the collision happens at the exact same "time" parameter, which is not required for geometric intersection.
\end{itemize}

\section{Part 4: Cheatsheet}

\begin{tcolorbox}[title=Essential Formulas]
\begin{itemize}
    \item Line: $\mathbf{r} = \mathbf{r}_0 + t\mathbf{v}$
    \item Plane: $\mathbf{n} \cdot (\mathbf{r} - \mathbf{r}_0) = 0$
    \item Normal from 3 Points: $\mathbf{n} = \vec{AB} \times \vec{AC}$
    \item Angle: $\cos \theta = \frac{|\mathbf{n}_1 \cdot \mathbf{n}_2|}{|\mathbf{n}_1||\mathbf{n}_2|}$
\end{itemize}
\end{tcolorbox}

\textbf{Quick Tips:}
\begin{itemize}
    \item \textbf{Line needs Direction ($\mathbf{v}$). Plane needs Normal ($\mathbf{n}$).} If you have the wrong one, use a Cross Product or geometric relation to swap.
    \item \textbf{Parallel Planes?} They share the same normal vector.
    \item \textbf{Line in Plane?} The line's direction $\mathbf{v}$ is perpendicular to the plane's normal $\mathbf{n}$ ($\mathbf{v} \cdot \mathbf{n} = 0$).
\end{itemize}

\textbf{Common Pitfalls:}
\begin{itemize}
    \item \textbf{Sign Errors in Symmetric Form:} Remember $\frac{x-x_0}{a}$. If you see $\frac{4-x}{2}$, rewrite it as $\frac{x-4}{-2}$. The coefficient of $x$ must be positive 1.
    \item \textbf{Obtuse Angles:} The angle between planes is usually defined as acute ($0 \le \theta \le 90^\circ$). If your dot product is negative, take the absolute value.
\end{itemize}

\newpage

\section{Part 5: Conceptual Synthesis}
\textbf{Thematic Connection:} The core theme is \textbf{Linearization}. We are describing flat, linear subspaces in $\mathbb{R}^3$. This is the precursor to Multivariable Calculus, where we will approximate curved surfaces using "Tangent Planes" (linear approximations). Just as the derivative gives a tangent line in 2D, the gradient will give a tangent plane in 3D.

\textbf{Forward Links:}
\begin{itemize}
    \item \textbf{Partial Derivatives:} You will use plane equations to visualize tangent planes to surfaces $z=f(x,y)$.
    \item \textbf{Lagrange Multipliers:} The condition $\nabla f = \lambda \nabla g$ implies parallel normal vectors, a concept grounded in this chapter.
\end{itemize}

\section{Part 6: Real-World Application (Finance Focus)}

\subsection{Scenario: Portfolio Optimization}
In quantitative finance, a portfolio of $n$ assets is a vector $\mathbf{w} = \langle w_1, w_2, \dots, w_n \rangle$. The constraint that fully invests the budget is $w_1 + \dots + w_n = 1$. In $\mathbb{R}^3$ (3 assets), this is the equation of a plane. The \textit{Efficient Frontier} involves maximizing return for risk, often visualized as a geometric problem involving the intersection of this "budget plane" with quadric surfaces (ellipsoids) representing constant risk.

\subsection{Model Setup}
\textbf{Problem:} An investor allocates wealth into 3 stocks. The weights must sum to 1. The expected return is fixed at 10\%.
\textbf{Variables:} $w_1, w_2, w_3$ (weights).
\textbf{Equations:}
1. Budget Plane: $w_1 + w_2 + w_3 = 1$ (Normal $\langle 1,1,1 \rangle$)
2. Return Plane: $r_1 w_1 + r_2 w_2 + r_3 w_3 = 0.10$ (Normal $\langle r_1, r_2, r_3 \rangle$)
\textbf{Solution:} The valid portfolios lie on the \textbf{line of intersection} of these two planes.

\section{Part 7: Untested Concepts}
\textbf{Distance from Point to Plane:}
Not covered in the homework, but standard.
Formula: $D = \frac{|ax_1 + by_1 + cz_1 + d|}{\sqrt{a^2+b^2+c^2}}$.
\textit{Example:} Distance from $(0,0,0)$ to $x+y+z=1$ is $\frac{|-1|}{\sqrt{3}} = 1/\sqrt{3}$.

\section{Part 8: Advanced Diagnostic - "Find the Flaw"}

\textbf{Problem 1: Symmetric Equations}
\textit{Student Work:} Line through $(1, 2, 3)$ with direction $\langle -2, 4, 1 \rangle$.
Eq: $\frac{1-x}{2} = \frac{y-2}{4} = z-3$.
\textbf{Error:} Incorrect variable order in first term.
\textbf{Correction:} $\frac{1-x}{2}$ is $\frac{-(x-1)}{2} = \frac{x-1}{-2}$. Correct form: $\frac{x-1}{-2} = \dots$

\textbf{Problem 2: Plane from Line}
\textit{Student Work:} Plane contains line $x=t, y=t, z=t$ and origin.
Vector $\mathbf{v}=\langle 1,1,1 \rangle$. Point $(0,0,0)$.
Conclusion: "Normal is $\langle 1,1,1 \rangle$".
\textbf{Error:} The student assumed the line direction is the normal.
\textbf{Correction:} The line is \textit{in} the plane. You need a second vector to cross with $\mathbf{v}$ to find $\mathbf{n}$. Since line passes through origin, you need a distinct point to form a second vector.

\textbf{Problem 3: Intersection}
\textit{Student Work:} Intersection of $x+y+z=1$ and $x+y+z=2$.
Solving... $1=2$. Error in algebra.
\textbf{Error:} Fundamental conceptual error.
\textbf{Correction:} The planes are parallel (same normal $\langle 1,1,1 \rangle$). They never intersect. Solution is the empty set.

\textbf{Problem 4: Angle Calculation}
\textit{Student Work:} Planes with normals $\langle 1,0,0 \rangle$ and $\langle -1, 1, 0 \rangle$.
Dot product $-1$. $\cos \theta = \frac{-1}{\sqrt{1}\sqrt{2}}$. $\theta = 135^\circ$.
\textbf{Error:} Angle between planes is defined as the acute angle.
\textbf{Correction:} Use absolute value of dot product. $\cos \theta = \frac{1}{\sqrt{2}} \implies 45^\circ$.

\textbf{Problem 5: Line Segment}
\textit{Student Work:} Parametrize segment from $A$ to $B$. $\mathbf{r}(t) = A + t(B-A)$.
"This represents the segment for all real $t$."
\textbf{Error:} Incorrect bounds.
\textbf{Correction:} The segment exists only for $0 \le t \le 1$. $t>1$ extends past $B$.

\end{document}